\section{The exact reach of convergent Cayley symmetry}
\label{cayley:sec:reach}

The existing $S_4$ proof uses the involution
$\phi(t)=(1-t)/(1+t)$ together with double shuffle.  Before attempting a
larger combined relation matrix, one can answer an exact structural
question: how much does this convergent involution prove by itself, and
how much stronger does it become after all shuffle products of its
relations are admitted?  These are two different presentations.  Both
can be computed in every weight.

The resulting dimensions below are dimensions of explicitly defined
formal word quotients.  They are not dimensions of evaluated periods.
The theorem includes neither stuffle nor distribution relations, so it
does not settle the manuscript's conjectured depth-two formula for $S_6$.

\subsection{Positive forms and two commuting involutions}

Use the five formal letters
\[
 z=\frac{dt}{t},\qquad c=\frac{dt}{1-t},\qquad
 a=\frac{i\,dt}{1-it},\qquad
 b=-\frac{dt}{1+t},\qquad
 \bar a=-\frac{i\,dt}{1+it}.
\]
Letters stand both for these forms and for a basis of a rational vector
space $V$.  Let $H(w)$ denote their outer-letter-first iterated integral
from zero to one.  An admissible nonempty word has first letter different
from $c$ and last letter different from $z$.  Each such integral is
absolutely convergent: the last-letter condition controls zero, and the
first-letter condition controls one.  The other endpoint singularities
are logarithmic and integrable.

Define the letter map $\tau=-\phi^*$ by
\begin{equation}\label{cayley:eq:tau}
\begin{array}{c|ccccc}
v&z&c&a&b&\bar a\\ \hline
\tau(v)&c-b&z+b&b-\bar a&b&b-a .
\end{array}
\end{equation}
Put
\[
 C(v_1\cdots v_n)=\tau(v_n)\cdots\tau(v_1),
 \qquad
 K(z)=z,\quad K(c)=c,\quad K(b)=b,\quad K(a)=\bar a,\quad K(\bar a)=a.
\]
Here products inside a word are concatenations; linear combinations are
expanded multilinearly.  Directly from the table,
\[
 C^2=K^2=1,\qquad CK=KC.
\]
Moreover $C$ preserves admissibility.  Indeed $\tau$ maps the span of
$\{z,a,b,\bar a\}$ to the span of $\{c,a,b,\bar a\}$ and maps the latter
span back to the former.

Changing variables and reversing the integration simplex gives
\begin{equation}\label{cayley:eq:period-law}
 H(w)=H(Cw)
\end{equation}
for every admissible word.  The sign from reversing the integration path
is included in the definition $\tau=-\phi^*$.  Thus there is no omitted
weight-dependent sign in \eqref{cayley:eq:period-law}.  Conjugation gives
$H(Kw)=\overline{H(w)}$.

Let $\mathcal A_n$ be the rational span of admissible words of length $n$,
and put $\mathcal A_0=\mathbb Q$.  Superscripts $+$ and $-$ on these spaces
will mean the $+1$ and $-1$ eigenspaces of $K$.  The minus space encodes
imaginary parts; choosing one word from each nontrivial conjugate pair is
an equivalent coordinate convention.

\subsection{All-weight rank of the unmultiplied Cayley rows}

\begin{theorem}[Raw convergent Cayley rank]\label{cayley:thm:raw}
For $n\ge2$,
\[
 d_n:=\dim_{\mathbb Q}\mathcal A_n^-
 =8\cdot5^{n-2}-2\cdot3^{n-2}.
\]
The span of the imaginary projections of all weight-$n$ rows $w-Cw$
has rank
\begin{equation}\label{cayley:eq:raw-rank}
 r_n=
 \begin{cases}
 4\cdot5^{n-2}-3^{n-2},&n\text{ even},\\[2mm]
 4\cdot5^{n-2}-3^{n-2}-2\cdot5^{(n-3)/2},&n\text{ odd}.
 \end{cases}
\end{equation}
At $n=1$, the imaginary space has dimension one and the rank is zero.
\end{theorem}

\begin{proof}
For $n\ge2$, write
\[
 U=\operatorname{span}\{z,a,b,\bar a\},\qquad
 E=\operatorname{span}\{c,a,b,\bar a\}.
\]
Then $\mathcal A_n=U\otimes V^{\otimes(n-2)}\otimes E$.
The dimensions of $U,E,V$ are $4,4,5$ and their conjugation traces
are $2,2,3$.  Therefore
\[
 \dim\mathcal A_n=16\cdot5^{n-2},\qquad
 \operatorname{tr}(K\mid\mathcal A_n)=4\cdot3^{n-2},
\]
which proves the formula for $d_n$.

It remains to compute traces of reversal.  For maps $T:U\to E$ and
$S:E\to U$, the map $u\otimes e\mapsto Se\otimes Tu$ has trace
$\operatorname{tr}(ST)$.  Applied to an endpoint pair of $C$, this
trace is four, because $\tau^2=1$.  Each paired pair of internal letters
contributes five.  The same facts hold for $CK$, because $(\tau K)^2=1$.
Finally the table gives
\[
 \operatorname{tr}(\tau\mid V)=1,\qquad
 \operatorname{tr}(\tau K\mid V)=-1.
\]
Thus, for $m\ge1$,
\[
\begin{array}{c|cc}
n&\operatorname{tr}C&\operatorname{tr}CK\\ \hline
2m&4\cdot5^{m-1}&4\cdot5^{m-1}\\
2m+1&4\cdot5^{m-1}&-4\cdot5^{m-1}.
\end{array}
\]
Because $C$ and $K$ commute,
\[
 \operatorname{tr}(C\mid\mathcal A_n^-)
 =\tfrac12\bigl(\operatorname{tr}C-\operatorname{tr}CK\bigr).
\]
The image of $1-C$ is precisely the negative eigenspace of $C$.
Its rank on $\mathcal A_n^-$ is consequently
$\tfrac12(d_n-\operatorname{tr}(C\mid\mathcal A_n^-))$.
Substitution proves \eqref{cayley:eq:raw-rank}.  Weight one is checked
directly: $K$ negates $a-\bar a$ and $C$ fixes it.
\end{proof}

In particular, the raw weight-seven Cayley matrix has $24\,514$
imaginary coordinates and rank $12\,207$.  There remain $12\,307$
coordinates in that particular quotient.  This count is exact even
though writing and row-reducing the full matrix is unnecessary.

\subsection{Closing the relations under shuffle multiplication}

The raw row span is not the same as the shuffle ideal generated by the
rows.  For example, a product of two $C$-negative elements is
$C$-positive; it vanishes in the ideal quotient but need not belong to
the raw negative eigenspace.

Let
\[
 \mathcal A=\bigoplus_{n\ge0}\mathcal A_n
\]
with the shuffle product.  Both $C$ and $K$ are graded algebra
automorphisms: reversal preserves a shuffle, as does applying a linear
letter map.  Define the homogeneous ideal and quotient
\[
 \mathcal J=\langle f-Cf:f\in\mathcal A\rangle_{\mathrm{shuffle}},
 \qquad \mathcal B=\mathcal A/\mathcal J.
\]
Thus every shuffle multiple of every convergent Cayley row is included.
The ideal is stable under $K$.

Let
\[
 h_n=\dim\mathcal B_n,\qquad
 g_n=\operatorname{tr}(K\mid\mathcal B_n),\qquad
 H(t)=\sum_{n\ge0}h_nt^n,\qquad G(t)=\sum_{n\ge0}g_nt^n.
\]
The letter $g_n$ in this subsection is a character trace, not a
Gaussian double $g_{a,b}$.

\begin{theorem}[The full Cayley shuffle-ideal quotient]
\label{cayley:thm:ideal}
The algebra $\mathcal B$ is a polynomial algebra on homogeneous
generators.  Its Hilbert and conjugation-character series are the unique
formal power series with constant term one satisfying
\begin{equation}\label{cayley:eq:hilbert}
 \frac{H(t)^2}{H(t^2)}=\frac{1-t}{1-5t},
 \qquad
 \frac{G(t)^2}{H(t^2)}
   =\frac{(1-t)^2}{(1+t)(1-3t)}.
\end{equation}
In particular,
\[
 \dim\mathcal B_n^-=\frac{h_n-g_n}{2},
 \qquad
 \dim\mathcal B_n^+=\frac{h_n+g_n}{2}.
\]
An explicit product form is
\begin{equation}\label{cayley:eq:dyadic-product}
 H(t)=\prod_{j\ge0}
 \left(\frac{1-t^{2^j}}{1-5t^{2^j}}\right)^{1/2^{j+1}},
\end{equation}
where each root is expanded with constant term one.
\end{theorem}

\begin{proof}
We supply the algebraic details, including the endpoint adjustment.
Let $\widetilde{\mathcal A}$ be the full shuffle algebra on $V$,
including words that are not admissible.  This is used only as a formal
algebra in the argument.

The shuffle polynomial theorem says that the Lyndon words freely
generate $\widetilde{\mathcal A}$ over $\mathbb Q$
\cite{cayley:radford,cayley:melancon-reutenauer}.
Use the alphabet order $z<a<b<\bar a<c$.  Every Lyndon word of length
at least two is admissible: it cannot begin with the greatest letter
$c$, or end with the least letter $z$.  Conversely, the nonincreasing
Lyndon factorization of an admissible word contains neither a
one-letter factor $c$ at the beginning nor a one-letter factor $z$
at the end.  The usual triangular shuffle factorization therefore
gives
\begin{equation}\label{cayley:eq:polynomial-endpoints}
 \widetilde{\mathcal A}\simeq\mathcal A[z,c]
\end{equation}
as graded polynomial algebras.  In particular $\mathcal A$ itself is
polynomial, and its indecomposables agree with those of
$\widetilde{\mathcal A}$ in every degree at least two.

For a connected graded algebra write
$Q(\mathcal A)=\mathcal A_{>0}/(\mathcal A_{>0})^2$.
Since $C$ and $K$ are commuting involutions over $\mathbb Q$, choose
a homogeneous simultaneous eigenbasis of $Q(\mathcal A)$.
Averaging arbitrary homogeneous lifts with the four character
projectors of the group generated by $C,K$ produces simultaneous
eigenvector lifts in $\mathcal A$.  These lifts freely generate
$\mathcal A$: replacing a graded polynomial generating set by lifts
of another basis of its indecomposables is a degree-triangular
invertible change of generators.

For this generating set, $\mathcal J$ is exactly the ideal generated
by the $C$-negative generators.  One inclusion follows from
$x-Cx=2x$ for each such generator.  For the other, after all
$C$-negative generators are set to zero, $C$ acts as the identity
on every remaining polynomial.  Thus every $f-Cf$ lies in that
ideal.  Consequently $\mathcal B$ is the polynomial algebra on
the $C$-positive generators, retaining their $K$-parities.

We compute the required generator characters without constructing
the generators.  The antipode of the full shuffle Hopf algebra is
$S(w)=(-1)^{|w|}\operatorname{rev}(w)$, and every connected graded
Hopf antipode acts as minus the identity on indecomposables.
Hence reversal acts on degree-$n$ indecomposables as
$(-1)^{n-1}$.  Put $\theta=-\tau$, applied letterwise.  On
indecomposables of every positive degree we therefore have
\[
 C=-\theta.
\]
In particular, the $C$-positive generator space is the
$\theta$-negative generator space, with the same $K$-characters.

Temporarily work in the full shuffle algebra.  Let $P(t)$ be the
Hilbert series of the polynomial algebra on its $C$-positive
generators, and $Q(t)$ the corresponding $K$-trace series.
The letter traces are
\[
\begin{array}{c|rrrr}
\sigma&1&\theta&K&K\theta\\ \hline
\operatorname{tr}(\sigma\mid V)&5&-1&3&1 .
\end{array}
\]
Since the degree-$n$ word space is $V^{\otimes n}$, the four
full-shuffle character series are respectively
\[
 \frac1{1-5t},\qquad\frac1{1+t},\qquad
 \frac1{1-3t},\qquad\frac1{1-t}.
\]

A $\theta$-negative polynomial generator of degree $n$ contributes
$(1+t^n)/(1-t^n)$ to the ratio of the identity-character series
to the $\theta$-character series.  Taking the product over all
such generators gives
\[
 \frac{P(t)^2}{P(t^2)}=\frac{1+t}{1-5t}.
\]
If the same generator has $K$-sign $\varepsilon\in\{1,-1\}$,
its contribution to the ratio of the $K$-character to the
$K\theta$-character is
$(1+\varepsilon t^n)/(1-\varepsilon t^n)$.
Since $\varepsilon^2=1$, this gives
\[
 \frac{Q(t)^2}{P(t^2)}=\frac{1-t}{1-3t}.
\]
These character products are valid for an infinite generating
set, since only finitely many generators contribute to any fixed
weight.

Finally, \eqref{cayley:eq:polynomial-endpoints} removes two
degree-one indecomposables and none in higher degree.
On their quotient space, $C$ interchanges $z$ and $c$ modulo
$\mathcal A_1$, while $K$ is the identity.  Exactly one removed
generator is $C$-positive, and it is $K$-positive.  Therefore
\[
 H(t)=(1-t)P(t),\qquad G(t)=(1-t)Q(t).
\]
Substituting these two relations proves
\eqref{cayley:eq:hilbert}.  Coefficients are determined recursively
because the coefficient of the new $h_n$ or $g_n$ in its square
is two.  Iterating the first equation proves
\eqref{cayley:eq:dyadic-product} in the $t$-adic topology.
\end{proof}

The first coefficients and their interpretation are
\[
\begin{array}{c|rrrrrrr}
n&1&2&3&4&5&6&7\\ \hline
\dim\mathcal A_n^-&1&6&34&182&946&4838&24514\\
r_n&0&3&15&91&463&2419&12207\\
h_n&2&9&36&162&720&3312&15336\\
g_n&0&3&4&18&36&124&300\\
\dim\mathcal B_n^-&1&3&16&72&342&1594&7518\\
\dim\mathcal J_n^-&0&3&18&110&604&3244&16996
\end{array}
\]
At weight seven, shuffle multiplication supplies $4789$ additional
imaginary relation directions beyond the raw Cayley row span.
The remaining $7518$ coordinates are independent only in this
specified formal quotient.  Adding double shuffle can reduce it
further; the dimension of an intersection of relation spaces is not
obtained by adding their separate ranks.

\subsection{Asymptotic size and the conjugation imbalance}

\begin{corollary}[Two distinct exponential scales]
\label{cayley:cor:asymptotic}
As $n\to\infty$,
\begin{align}
 h_n&=
 \sqrt{\frac{4H(1/25)}{5\pi}}\,
 \frac{5^n}{\sqrt n}\left(1+O(n^{-1})\right),\\
 g_n&=
 \sqrt{\frac{H(1/9)}{3\pi}}\,
 \frac{3^n}{\sqrt n}\left(1+O(n^{-1})\right),\\
 \dim\mathcal B_n^-&=
 \sqrt{\frac{H(1/25)}{5\pi}}\,
 \frac{5^n}{\sqrt n}\left(1+O(n^{-1})\right).
\end{align}
Both of the first two expansions admit all orders in $1/n$.
More precisely, the relative difference between the real and
imaginary dimensions satisfies
\[
 \frac{\dim\mathcal B_n^+-\dim\mathcal B_n^-}
      {\dim\mathcal B_n^++\dim\mathcal B_n^-}
 =
 \sqrt{\frac{5H(1/9)}{12H(1/25)}}\,
 \left(\frac35\right)^n\left(1+O(n^{-1})\right).
\]
\end{corollary}

\begin{proof}
The product \eqref{cayley:eq:dyadic-product} is convergent and
nonvanishing in $|t|<1/5$.  The functional equations give
\[
 H(t)=F(t)(1-5t)^{-1/2},\qquad
 F(t)=\sqrt{(1-t)H(t^2)},
\]
and
\[
 G(t)=E(t)(1-3t)^{-1/2},\qquad
 E(t)=(1-t)\sqrt{\frac{H(t^2)}{1+t}}.
\]
The factors $F,E$ are analytic on the disk $|t|<1/\sqrt5$;
the roots are the analytic roots taking value one at zero.
Consequently the only dominant singularity for $H$ is $1/5$,
and that for $G$ is $1/3$.  The standard binomial coefficient
asymptotic
$[t^n](1-\lambda t)^{-1/2}
=\lambda^n(\pi n)^{-1/2}(1+O(n^{-1}))$,
together with Taylor expansion of the analytic factor at the
singularity, yields the first two assertions.  The factor values are
\[
 F(1/5)=\sqrt{\tfrac45 H(1/25)},\qquad
 E(1/3)=\sqrt{\tfrac13 H(1/9)}.
\]
The same Taylor expansion gives every further inverse-power
coefficient.  Finally
$\dim\mathcal B_n^\pm=(h_n\pm g_n)/2$ and division proves the ratio.
In particular, the $3^n$ character contribution is not absorbed into
an assertion that a finite $5^n/n$ expansion resolves exponentially
smaller terms: it has its own independently defined series $G$.
\end{proof}

\subsection{A completely explicit Cayley family containing \texorpdfstring{$S_6$}{S6}}

Retain
$S_p=\sum_{n\ge0}(-1)^nH_n/(2n+1)^p$.
The harmonic parity split and the positive-form convention give
\[
 S_p=\operatorname{Im}H\bigl(z^{p-1}a(a+\bar a)\bigr).
\]
This follows absolutely for $p>1$, and by a radial Abel limit for
$p=1$.

\begin{proposition}[An all-index pure-depth representation]
\label{cayley:prop:harmonic}
For every integer $p\ge1$,
\begin{equation}\label{cayley:eq:harmonic}
 S_p=\operatorname{Im}H\left((2ba-aa+a\bar a)(c-b)^{p-1}\right).
\end{equation}
All products and powers inside $H$ are concatenations.
The expanded right side contains exactly $3\cdot2^{p-1}$
admissible words, each of weight and series depth $p+1$.
In particular
\[
 \boxed{\quad
 S_6=\operatorname{Im}H\left((2ba-aa+a\bar a)(c-b)^5\right).
 \quad}
\]
\end{proposition}

\begin{proof}
Cayley transport gives
\[
 C\bigl(z^{p-1}a(a+\bar a)\bigr)
 =(2b-a-\bar a)(b-\bar a)(c-b)^{p-1}.
\]
Both $2b-a-\bar a$ and $c-b$ are fixed by conjugation.
Taking half the difference with the conjugate and then grouping
the three conjugate pairs gives exactly the imaginary part in
\eqref{cayley:eq:harmonic}.  The first letters in the resulting
three word families are $a$ or $b$, and all letters are nonzero
kernel letters.  Thus every expanded word is admissible and
represents an all-one colored multiple polylogarithm.
\end{proof}

This proposition provides a new explicit high-depth expression
and a compact formal certificate for the weight-seven mixed sum.
It does not turn the manuscript's proposed depth-two $S_6$
reduction into a theorem.  The latter still requires an exact
combination of additional relations.

For a real-integral check of conventions, the same identity yields
\[
 S_p=\frac1{(p-1)!}\int_0^1
 \log^{p-1}\!\left(\frac{1+u}{1-u}\right)
 \frac{\log\!\left((1+u)^2/(2(1+u^2))\right)}{1+u^2}\,du.
\]
Indeed the iterated integral of $(c-b)^{p-1}$ up to $u$ is the
displayed logarithmic power divided by $(p-1)!$.
Integrating the two outer letters in
$2ba-aa+a\bar a$, and taking imaginary parts, produces the
remaining kernel.  The substitution $u=(1-x)/(1+x)$ returns
the original Mellin representation of $S_p$.

There is also a direct verification of this real integral, independent
of the word expansion.  Start from the already convergent identity
\[
 S_p=-\frac1{(p-1)!}\int_0^1
 (-\log x)^{p-1}\frac{\log(1+x^2)}{1+x^2}\,dx.
\]
With $x=(1-u)/(1+u)$ the endpoint pair $(0,1)$ becomes $(1,0)$, and
\[
 -\log x=\log\frac{1+u}{1-u},\qquad
 \frac{dx}{1+x^2}=-\frac{du}{1+u^2},\qquad
 \log(1+x^2)=-\log\frac{(1+u)^2}{2(1+u^2)}.
\]
Substitution, including the reversed integration limits, gives the
displayed formula.  Near $u=1$ the new logarithmic kernel vanishes
quadratically in $1-u$, so every fixed logarithmic power remains
integrable.  At $u=0$ the integrand is bounded (and has a zero of
order $p-1$ when $p>1$).

\subsection{Finite verification and the remaining mixed reduction}

The supplied program \texttt{check\_cayley\_ranks.py} is independent
of the older report's word implementation.  It builds all admissible
words, expands the displayed letter transformation, checks $C^2=1$
and $CK=KC$, and computes rational ranks by integer elimination
with gcd normalization.  It also constructs every shuffle lift
$H(v)H(w-Cw)$ with total weight at most four.

The exact matrix results are
\[
\begin{array}{c|rrrr}
n&1&2&3&4\\ \hline
\dim\mathcal A_n&3&16&80&400\\
\text{raw full rank}&1&6&38&190\\
\text{all lifted full rank}&1&7&44&238\\
\text{raw imaginary rank}&0&3&15&91\\
\text{all lifted imaginary rank}&0&3&18&110
\end{array}
\]
The weight-four lifted calculation contains all $1136$ prescribed
row instances.  These are exact rational checks, not modular rank
estimates.  They audit the all-weight proofs; the proofs do not
depend on extrapolating the table.

The separate program \texttt{verify\_harmonic\_cayley.py} rebuilds
\eqref{cayley:eq:harmonic} for $1\le p\le8$ and obtains zero exact
residual in every case.  The file
\texttt{S6\_cayley\_depth\_seven.json} lists all $96$ terms for
$p=6$.  The verifier imports no numerical polylogarithm evaluator.

An exploratory weight-seven double-shuffle computation generated
$51\,244$ normalized rows, in addition to a selected Cayley row.
The sparse search was stopped after substantial fill-in; it
produced neither an exact certificate for the conjectured $S_6$
formula nor a separating functional.  No membership or
nonmembership claim follows from that unfinished computation.
The exact quotient and rank theorems above concern the fully
specified Cayley-only presentations and remain unaffected.
